\section{Globals class}

The {\tt globals} class is used to refer to all state variables as one
object, for the puposes of things like checkpoint restart. The object
may be a subrange of cells, for example the cells that are represented
on a particular processor. Methods exist to extract or replace
subranges of cells. The index operator allows components of the arrays
to be extracted symbolically. So if {\tt density}$\equiv${\tt
  global\_vars.iarrays[0]}, then {\tt g[density]} returns {\tt g.iarrays[0]}.

\begin{verbatim}
//A class for storing the global variables used in the model
class global
{
public:
  //origin cell and number of cells on this processor 
  int ocell, ncells, tstep; 

  // dimensions array and number of species in each cell
  iarray dims, nsp, *coords;  

  // list of iarrays
  iarray *iarrays;

  //list of arrays
  array *arrays;

  //list of sparse_mats
  sparse_mat *sparse_mats;

  //lengths of above lists
  int niarr, narr, nspars;

  global(); 
  /* i=no. of iarrays, a=no. of arrays and s=no. of sparse_mats */
  global(int i, int a, int s);
  global(global&); 
  global operator=(global);
  ~global();

  /* get local equivalents of a global variable */
  iarray& operator[](iarray& var); 
  array& operator[](array& var);
  sparse_mat& operator[](sparse_mat& var);
  

  /* return just the global variables belonging to cell, or a number
      of cells cell, cell+1, ... cell+nc-1 */
  global get_cell_vars(int cell, int nc=1);
 
  //replace cells variables by those contained in vars
  void put_cell_vars(global vars);
  //add a cell into this global
  void append(global vars);

  //pack up the global data into a contiguous package
  void packup(glue& buffer);
  //unpack contiguous representation of global data
  void unpack(glue& buffer);
};

\end{verbatim}


The {\tt glue} class is a helper class that represents the global
variables contiguously. It has four important members; {\tt char
  *data} and {\tt int size}, which point to the contiguous
representation, and the number of bytes within that representation;
and the methods \verb|<<| and \verb|>>|, which appends/extracts an
object into/from the representation. The {\tt packup} and {\tt unpack}
methods within the globals class convert between the glue
representation and the globals representation.

\begin{verbatim}
//A Class for grouping a whole bunch of binary things together
class glue
{
public:
  char *data;
  int size;
  glue(int sz);
  ~glue()
  void operator<<(int);
  void operator<<(double);
  void operator<<(iarray);
  void operator<<(array);
  void operator<<(sparse_mat);
  void operator>>(int&);
  void operator>>(double&);
  void operator>>(iarray&);
  void operator>>(array&);
  void operator>>(sparse_mat&);
};

\begin{verbatim}
//Assemble a full set of global variables by combining data on each processor
void get_all_globals(global&);
\end{verbatim}

collect the global variables on each processor onto processor 0.
